\documentclass[11pt, a4paper]{amsart}
\usepackage{amsmath, amssymb, amsthm, amsfonts}
\usepackage{mathtools}
\usepackage[margin=1in]{geometry}
\usepackage{hyperref}

\newtheorem{theorem}{Theorem}[section]
\newtheorem{lemma}[theorem]{Lemma}
\newtheorem{proposition}[theorem]{Proposition}
\newtheorem{corollary}[theorem]{Corollary}
\newtheorem{fact}[theorem]{Standard Fact}
\theoremstyle{definition}
\newtheorem{definition}[theorem]{Definition}
\newtheorem{remark}[theorem]{Remark}
\newtheorem{notation}[theorem]{Notation}
\newtheorem{hypothesis}[theorem]{Hypothesis}
\newtheorem{axiom}[theorem]{Axiom}

\newcommand{\R}{\mathbb{R}}
\newcommand{\C}{\mathbb{C}}
\newcommand{\Q}{\mathbb{Q}}
\newcommand{\Z}{\mathbb{Z}}
\newcommand{\HH}{\mathbb{H}}
\newcommand{\rk}{\operatorname{rank}}
\newcommand{\Gal}{\operatorname{Gal}}
\newcommand{\Frac}{\operatorname{Frac}}
\newcommand{\Hmin}{\mathbf{H}_{\min}}
\newcommand{\Hor}{\mathbf{H}_{\mathrm{or}}}
\newcommand{\Hdist}{\mathbf{H}_{\mathrm{dist}}}
\newcommand{\Hquot}{\mathbf{H}_{\mathrm{quot}}}
\newcommand{\Ophi}{\mathbb{Z}[\varphi]}
\newcommand{\Mcox}{M_{\mathrm{Cox}}}
\newcommand{\Mtw}{M_{12}}
\newcommand{\Wph}{W_{\mathrm{ph}}}
\newcommand{\Nslot}{N_{\mathrm{slot}}}

\title{Phason--Coxeter Alignment:\\
  Coxeter-Plane Resolution of the 12D Phason Complement\\
  and Upper-Exponent Winding}
\author{}
\date{}

\begin{document}
\maketitle

\begin{abstract}
This is the second entry of the $\alpha$-chain in the
cascade-correspondence programme, building strictly on
$\alpha$-1 (\emph{The 12D Closure Theorem}) and on classical
Lie / Coxeter / icosian / McKay theory. It does \emph{not}
depend on the infrastructure papers P1--P7.

We carry out two steps. First (\emph{Theorem~A}, $T_A$,
Theorem~\ref{thm:TA}), \emph{relative to a fixed Coxeter
element $c \in W(E_8)$} and under $\alpha$-1's Hypothesis
$\Hmin$ together with Hypothesis $\Hor$
(Hypothesis~\ref{hyp:Horient}, the $\sigma$-twist
orientation convention on Coxeter planes), we show that
the $\alpha$-1 phason complement $\Mtw \cong \Ophi^{2}$
admits a Coxeter-resolved bookkeeping module $\Mcox(c)
\cong \Ophi^{4}$ with one $\Ophi$-line per
$c$-invariant Coxeter plane $P_k(c)$, canonical relative
to $c$ up to signed permutation $S_4 \ltimes (\Z/2)^4$
of plane labels (across distinct choices of $c$, only
existence-up-to-conjugacy is asserted). The bookkeeping
equivalence $\Mtw \leftrightarrow \Mcox(c)$ is proved at
the cut-and-project parameter level, not the
ambient-lattice level: $L_{16}(c) := E_8 \oplus \Mcox(c)$
is \emph{not} a new closure theorem, only a four-plane
resolution of the same underlying internal geometry of
$\alpha$-1's $L_{12}$.

Second (\emph{Theorem~B}, $T_B$, Theorems~\ref{thm:TB-88}
and~\ref{thm:TB-87}), under $T_A$ plus an explicit
$\sigma$-twist orientation hypothesis $\Hor$
(Hypothesis~\ref{hyp:Horient}), each plane-wise phason
winding integer $w_k$ equals the \emph{upper} exponent
$m_k^+$ in the conjugate exponent pair attached to $P_k$.
The total phason winding is
\[
  \Wph \;=\; \sum_{k=1}^{4} m_k^+
        \;=\; 17 + 19 + 23 + 29
        \;=\; 88,
\]
proved as Theorem~\ref{thm:TB-88} \emph{conditionally on
$\Hmin + \Hor + \Hdist$}, where $\Hdist$
(Hypothesis~\ref{hyp:Hdist}) is a distribution hypothesis
ensuring the four plane-wise images of the reconciliation
map $\Psi$ together generate the phason internal image.
The slot-count statement $\Nslot = 87$ requires the
additional named hypothesis $\Hquot$
(Hypothesis~\ref{hyp:Hquot}), which fixes a one-dimensional
global-phase quotient; this clause is stated in
Theorem~\ref{thm:TB-87}, conditionally on $\Hmin + \Hor +
\Hdist + \Hquot$.

The paper makes no claim about $\alpha^{-1}$, the
$\pi/87$ correction, the $137 = 87 + 50$ decomposition,
the $600$-cell Laplacian spectrum, photons, microtubules,
or Standard-Model phenomenology. Those belong to $\alpha$-3
and to downstream papers. We make no claim that
$L_{16}$ replaces $L_{12}$ as the canonical upward closure.
\end{abstract}

\tableofcontents

%=====================================================================
\section{Introduction and Convention Reconciliation}
\label{sec:intro}
%=====================================================================

\subsection{What this paper is}

This paper is the second entry of the $\alpha$-chain. It
takes the $12$D closure structure of $\alpha$-1
\cite{Cascade12D} as input, together with classical Lie /
Coxeter / icosian / McKay theory, and derives two
structural results about how the $\alpha$-1 phason complement
sits inside the Coxeter-plane decomposition of $E_8$.

The position of the $\alpha$-chain in the cascade-correspondence
programme is described in $\alpha$-1's introduction; we recall
only that the $\alpha$-chain is \emph{parallel} to the
infrastructure stack P1--P7 and shares with it only the
classical $E_8$ / icosian core. No infrastructure paper is
used as a proof source here.

\subsection{What this paper proves}

\textbf{Unconditional layer.}
\begin{description}
\item[\textbf{Lemma~\ref{lem:E8-data}.}] Classical $E_8$
  Coxeter data: Coxeter number $h = 30$, exponents
  $\{1, 7, 11, 13, 17, 19, 23, 29\}$, conjugate pairing
  $m \leftrightarrow 30 - m$, four pairwise-orthogonal
  Coxeter planes $P_1, \ldots, P_4 \subset \R^8$.
\item[\textbf{Lemma~\ref{lem:cox-plane}.}] Kostant's Coxeter-plane
  decomposition: $\R^8 = P_1 \oplus P_2 \oplus P_3 \oplus P_4$
  pairwise-orthogonally, with $c$ acting as a rotation in each
  plane.
\item[\textbf{Lemma~\ref{lem:McKay}.}] McKay correspondence
  for $2I$ and $E_8$: classical ADE bridge linking the binary
  icosahedral group to the $E_8$ Dynkin diagram.
\end{description}

\textbf{Conditional on $\Hmin + \Hor + \Hdist$, relative to
a fixed Coxeter element $c \in W(E_8)$:}
\begin{description}
\item[\textbf{Theorem~\ref{thm:TA}}] (\textbf{T${}_A$}, the
  phason--Coxeter alignment): for the chosen $c$, the
  Coxeter-resolved phason module $\Mcox(c)$ decomposes as
  \[
    \Mcox(c) \;\cong\; \bigoplus_{k=1}^{4} \Ophi \cdot e_k,
  \]
  with one $\Ophi$-line per $c$-invariant Coxeter plane
  $P_k(c)$, canonical (relative to $c$) up to signed
  permutation $S_4 \ltimes (\Z/2)^4$ of the four plane
  labels. The reconciliation map $\Psi$ to the $\alpha$-1
  phason image is constructed explicitly plane-by-plane.
  Across distinct choices of $c$, only existence up to
  $W(E_8)$-conjugacy is asserted.
\item[\textbf{Theorem~\ref{thm:TB-88}.}] Total phason
  winding $\Wph = 88$. Conditionality flows through $T_A$
  (using $\Hmin$, $\Hor$, $\Hdist$) and through the
  plane-wise upper-exponent labelling (which uses $\Hor$
  directly).
\end{description}

\textbf{Conditional on $\Hmin + \Hor + \Hdist + \Hquot$:}
\begin{description}
\item[\textbf{Theorem~\ref{thm:TB-87}}] (\textbf{T${}_B$,
  slot count}): under all four hypotheses (with $\Hquot$
  fixing the one-dimensional global-phase quotient), the
  slot count is $\Nslot = \Wph - 1 = 87$.
\end{description}

\subsection{What this paper does NOT claim}

\begin{enumerate}
\item No claim that $T_A$ is unconditional. It depends on
  $\Hmin$ (via $\alpha$-1's $C_1$), $\Hor$, and $\Hdist$.
\item No claim that $\Wph = 88$ depends on $\Hquot$. The
  total winding count is established under $\Hmin + \Hor +
  \Hdist$, \emph{before} the one-dimensional quotient is
  taken.
\item No claim that $\Hor$, $\Hdist$, or $\Hquot$ is
  itself a theorem. All are named hypotheses with structural
  justifications (Remarks~\ref{rmk:Horient-status},
  \ref{rmk:Hdist-status}, \ref{rmk:Hquot-justif}); the
  slot-count clause $\Nslot = 87$ is stated conditionally on
  $\Hmin + \Hor + \Hdist + \Hquot$.
\item No claim about $\alpha^{-1} = 137 + \pi/87$. That is the
  subject of $\alpha$-3.
\item No claim that $137 = 87 + 50$. The $50$ count and the
  $137$ decomposition are not addressed here.
\item No claim about the $600$-cell rational-irrep Laplacian
  spectrum.
\item No claim about photons, microtubules, biology, the
  Standard Model, or any phenomenology.
\item No claim that $L_{16} = E_8 \oplus \Mcox$ is a new closure
  theorem. $L_{16}$ is a Coxeter-resolved bookkeeping object
  attached to the same underlying $L_{12}$ closure of
  $\alpha$-1; we make no statement that $L_{16}$ replaces
  $L_{12}$ as the canonical upward ambient.
\item No claim of dependence on P1--P7. Citations of the
  infrastructure stack are absent throughout.
\item No use of \texttt{papers/cascade-derivation/} working
  notes as theorem-grade proof sources. Those notes are author
  scratch material; their contents are converted to clean
  theorem statements with classical citations here.
\end{enumerate}

\subsection{Convention Reconciliation}
\label{sec:conv-recon}

A central technical point is that this paper uses a different
rank reading of the phason complement than $\alpha$-1
\cite{Cascade12D}. We fix the relationship explicitly here,
\emph{before} any theorem.

\paragraph{$\alpha$-1's reading.} Under $\alpha$-1's
Notation \texttt{not:rank-convention}, ``rank'' means $\rk_\Z$.
The closure ambient is $L_{12}$ with $\rk_\Z L_{12} = 12$, and
under Hypothesis $\Hmin$, the phason complement $\Mtw$ has
$\rk_\Z \Mtw = 4$, equivalently $\rk_\Ophi \Mtw = 2$:
$\Mtw \cong \Ophi^{2}$ as $\Ophi$-modules (this is $\alpha$-1's
Theorem $C_1$).

\paragraph{$\alpha$-2's resolved reading.} For the
Coxeter-plane decomposition, the load-bearing object is the
\emph{four-line Coxeter-resolved bookkeeping module}
\[
  \Mcox \;:=\; \bigoplus_{k=1}^{4} \Ophi \cdot e_k,
\]
free of $\Ophi$-rank $4$ (equivalently $\Z$-rank $8$), one
$\Ophi$-line per Coxeter plane. We adopt the
\emph{$\Ophi$-rank reading} for the bookkeeping module from
this point on: when we write ``$\Mcox \cong \Ophi^{4}$'' the
exponent is $\Ophi$-rank, not $\Z$-rank.

\paragraph{Bookkeeping equivalence.} The relationship between
$\Mtw$ and $\Mcox$ is \emph{not} a module isomorphism (their
$\Z$-ranks differ: $4$ vs $8$). Rather, $\Mcox$ is a
\emph{Coxeter-resolved parameter module} for the same
internal-image data carried by $\Mtw$ in the cut-and-project
package: each plane-wise component $\Ophi \cdot e_k$ encodes
the shift amplitude of the underlying $\sigma$-embedding of
$E_8$ along the Coxeter plane $P_k$. Under the conjunction
$\Hmin + \Hor + \Hdist$, the $\Ophi$-action structures of
$\Mtw$ and $\Mcox$ match in a precise sense made explicit in
Definition~\ref{def:reconciliation} and
Proposition~\ref{prop:reconciliation}. We do \emph{not} claim
that the formal sum $L_{16} := E_8 \oplus \Mcox$ is a 16D
closure ambient analogous to $\alpha$-1's $L_{12}$; that
would be a separate (unproven) closure theorem. $L_{16}$ is
used in this paper purely as bookkeeping notation for the
Coxeter-resolved phason data.

\paragraph{Notation discipline.} After this subsection,
``rank'' \emph{always} means $\rk_\Ophi$ (Coxeter-resolved
$\Ophi$-rank reading); every $\Z$-rank is marked
$\rk_\Z$ explicitly. When citing $\alpha$-1's $L_{12}$ and
$\Mtw$ we will retranslate using the bookkeeping equivalence
of Proposition~\ref{prop:reconciliation}.

\subsection{Source audit}

\textbf{Internal theorem inputs.}
\begin{itemize}
\item Axiom $\mathbf{F1}$ (the permeability fixed-point axiom)
  is a named axiom of the cascade programme whose statement is
  collected in \cite[\S F1]{CascadeFoundations}; we adopt it
  here as an axiom (not a theorem) and use only its statement,
  not any derivation from $\mathbf{F1}$ that \cite{CascadeFoundations}
  might contain. We rely on $\mathbf{F1}$ only in
  \S\ref{sec:convention}'s convention reconciliation for the
  $L_{12}$ / $L_{16}$ bookkeeping, where the axiom fixes the
  permeability balance.
\item $\alpha$-1 Theorems $C_1$, $C_2$, $C_3$
  \cite{Cascade12D}: $C_2$ may be cited unconditionally for
  the $4 + 4$ closure architecture; $C_1$ and $C_3$ may be
  cited only conditionally on $\Hmin$ (their statement
  hypothesis in $\alpha$-1).
\end{itemize}

\textbf{Classical mathematical inputs.}
\begin{itemize}
\item \emph{Coxeter / Lie theory.}
  \cite{Bourbaki46} for exponents and Coxeter elements of
  $E_8$ (Ch.~6 \S4 No.~12, Planche~VII for $E_8$);
  \cite{Humphreys} as a clean modern reference (\S 3.16--3.20
  for Coxeter elements and exponents).
\item \emph{Coxeter-plane decomposition.}
  \cite{Kostant59} for the principal $\mathfrak{sl}_2$
  decomposition into pairwise-orthogonal $c$-invariant
  $2$-planes (Theorem 6.7 and surrounding discussion).
\item \emph{Icosian / $E_8$ / $H_4$ / 600-cell.}
  \cite{ConwaySloane} for the icosian ring $I$ and the
  $E_8 = \{(x, \sigma(x)) : x \in I\}$ embedding (Ch.~8 \S2);
  \cite{MoodyPatera} for the $E_8 \to H_4$ cut-and-project
  (Thm~4.1, Prop.~4.2); \cite{ElserSloane} for the
  cut-and-project quasicrystal construction (\S2--3,
  pp.~6162--6165).
\item \emph{McKay correspondence.}
  \cite{McKay80} for the ADE / finite-subgroup-of-$SU(2)$
  correspondence linking $2I$ to $E_8$.
\end{itemize}

We do not cite P1--P7, the \texttt{cascade-derivation/} working
notes, scripts, notebooks, or downstream $\alpha$-3 / photon /
microtubule manuscripts as proof sources.

\subsection{Organization}

\S\ref{sec:prelims} states $\mathbf{F1}$ as an axiom (carried
over from $\alpha$-1), records the Galois data $\sigma$, and
imports the $\alpha$-1 Theorems used here.

\S\ref{sec:cox-data} proves Lemmas~\ref{lem:E8-data}
($E_8$ exponent / Coxeter-number data),
\ref{lem:cox-plane} (Kostant's plane decomposition), and
\ref{lem:McKay} (McKay correspondence).

\S\ref{sec:Mcox} introduces the Coxeter-resolved phason
module $\Mcox$ and the reconciliation map relating it to
$\alpha$-1's $\Mtw$ (Definition~\ref{def:reconciliation},
Proposition~\ref{prop:reconciliation}).

\S\ref{sec:TA} proves Theorem~\ref{thm:TA} ($T_A$:
phason--Coxeter alignment) under $\Hmin + \Hor + \Hdist$.

\S\ref{sec:sigma-twist} establishes the $\sigma$-twist
orientation reversal (Hypothesis~\ref{hyp:Horient}) and defines
plane-wise winding integers (Definition~\ref{def:winding}).

\S\ref{sec:TB-88} proves Theorem~\ref{thm:TB-88}
($\Wph = 88$) under $\Hmin + \Hor + \Hdist$.

\S\ref{sec:TB-87} introduces Hypothesis $\Hquot$
(one-dimensional global-phase quotient) and proves
Theorem~\ref{thm:TB-87} ($\Nslot = 87$ under $\Hmin + \Hor +
\Hdist + \Hquot$).

\S\ref{sec:audit} is the self-audit.
\S\ref{sec:handoff} is the citation contract.
\S\ref{sec:notation} is the notation audit.

%=====================================================================
\section{Preliminaries: $\alpha$-1 Imports and Galois Data}
\label{sec:prelims}
%=====================================================================

\begin{axiom}[$\mathbf{F1}$: the permeability fixed-point
  axiom, recap from $\alpha$-1]
\label{ax:F1}
There exists a positive real $\varphi$ satisfying
$\varphi = 1 + 1/\varphi$, equivalently $\varphi^2 = \varphi
+ 1$, whence $\varphi = (1 + \sqrt{5})/2$. We import this
axiom from \cite[\S F1]{CascadeFoundations} via $\alpha$-1
\cite[Axiom~ax:F1]{Cascade12D}.
\end{axiom}

\begin{notation}[Number-field data]
\label{not:K}
Let $K := \Q(\sqrt 5) = \Q(\varphi)$, $\Ophi := \Z[\varphi]
= \mathcal O_K$ (the ring of integers; see
\cite[Ch.~2, Thm.~1]{Marcus}), and $\sigma \in \Gal(K / \Q)$
the nontrivial Galois automorphism, $\sigma(\varphi) = 1 -
\varphi$. As proved in $\alpha$-1, $\Ophi$ is a PID
\cite[Lem.~lem:pid]{Cascade12D} and every finitely generated
torsion-free $\Ophi$-module is free
\cite[Lem.~lem:pid-free]{Cascade12D}.
\end{notation}

\begin{notation}[Rank convention for this paper]
\label{not:rank-conv}
\emph{From this point on}, ``rank'' and $\rk$ always mean
$\rk_\Ophi$ (the Coxeter-resolved $\Ophi$-rank reading
adopted in \S\ref{sec:conv-recon}). Whenever we refer to
$\Z$-rank, we write $\rk_\Z$ explicitly. The $\alpha$-1
ambient $L_{12}$ has $\rk_\Z L_{12} = 12$ and $\rk_\Ophi
L_{12} = 6$ (its $\Ophi$-rank reading); $\Mtw$ has
$\rk_\Z = 4$ and $\rk_\Ophi = 2$. The Coxeter-resolved
bookkeeping module $\Mcox$ has $\rk_\Ophi = 4$ and
$\rk_\Z = 8$.
\end{notation}

\begin{fact}[Imports from $\alpha$-1]
\label{fact:alpha1-imports}
We freely use the following from
$\alpha$-1 \cite{Cascade12D}.
\begin{enumerate}
\item[(a)] $\alpha$-1 Theorem $C_2$ (\emph{unconditional}):
  given the classical Elser--Sloane package
  (Definitions \texttt{def:cap}, \texttt{def:E8-iso} of
  $\alpha$-1) and the admissibility class fixed by
  axioms (A1)--(A4), no admissible upward ambient of
  $\Z$-rank $> 12$ adds new geometric content; the canonical
  closure has $\rk_\Z L_{12} = 12$.
\item[(b)] $\alpha$-1 Hypothesis $\Hmin$
  (\emph{conditional}): the phason coefficient ring of
  $\Mtw$ in $L_{12} = E_8 \oplus \Mtw$ is $\Ophi$.
\item[(c)] $\alpha$-1 Theorem $C_1$ (\emph{conditional on
  $\Hmin$}): under $\Hmin$, $\Mtw \cong \Ophi^{2}$ as an
  $\Ophi$-module (the phason complement is free of
  $\Ophi$-rank $2$, equivalently $\Z$-rank $4$).
\item[(d)] $\alpha$-1 Theorem $C_3$ (\emph{conditional on
  $\Hmin$}): the phason coefficient ring at $L_{12}$ closure
  coincides with the smallest subring of $\R$ containing
  $\Z$ and the $\mathbf{F1}$ fixed point $\varphi$.
\end{enumerate}
\end{fact}

%=====================================================================
\section{Coxeter-Plane Structure of $E_8$}
\label{sec:cox-data}
%=====================================================================

We set up the classical Coxeter-plane data needed for
$T_A$ and $T_B$. Each lemma in this section is a citation of
classical theory; we record the statements verbatim for use
downstream and pin the source.

\begin{lemma}[$E_8$ Coxeter data]
\label{lem:E8-data}
The Weyl group $W(E_8)$ acts on the reflection representation
$V := \R^8$ (the span of the $E_8$ root system). A Coxeter
element $c \in W(E_8)$ has order $h = 30$ (the Coxeter number).
The exponents of $E_8$ are
\[
  \{m_1, \ldots, m_8\}
  \;=\; \{1, 7, 11, 13, 17, 19, 23, 29\},
\]
each appearing with multiplicity one and satisfying the
conjugate pairing
\[
  m_k^- + m_k^+ \;=\; 30
  \qquad(k = 1, 2, 3, 4),
\]
with the four conjugate pairs
\[
  \{1, 29\}, \quad \{7, 23\}, \quad \{11, 19\}, \quad \{13, 17\}.
\]
The lower exponents are
$\{m_k^-\} = \{1, 7, 11, 13\}$ and the upper exponents are
$\{m_k^+\} = \{17, 19, 23, 29\}$, with
\[
  \sum_{k=1}^{4} m_k^+ \;=\; 17 + 19 + 23 + 29 \;=\; 88.
\]
\end{lemma}

\begin{proof}
The Coxeter number, the exponent list, and their multiplicities
for $E_8$ are classical. See
\cite[Ch.~6, \S4 No.~12, Planche~VII (the table for $E_8$)]{Bourbaki46}
or \cite[\S 3.18 (Coxeter number) and \S 3.19 (exponents)
together with the table in \S 3.7 for $E_8$]{Humphreys}.
The pairing $m \leftrightarrow h - m$ follows from
\cite[Ch.~5, \S6 No.~2, Cor.~2 (the symmetry of exponents
about $h/2$)]{Bourbaki46} or
\cite[Thm.~3.18 and Cor.~3.20]{Humphreys}.
The arithmetic $17 + 19 + 23 + 29 = 88$ is direct.
\end{proof}

\begin{lemma}[Kostant's Coxeter-plane decomposition]
\label{lem:cox-plane}
Let $c \in W(E_8)$ be a Coxeter element acting on $V = \R^8$.
There is an orthogonal direct-sum decomposition
\[
  V \;=\; P_1 \oplus P_2 \oplus P_3 \oplus P_4,
\]
where each $P_k$ is a real $2$-dimensional $c$-invariant
subspace, the four planes are pairwise orthogonal in the
$W(E_8)$-invariant inner product on $V$, and on each $P_k$,
$c$ acts as rotation by angle $2\pi m_k^- / 30$ (the lower
exponent in the conjugate pair attached to $P_k$).
\end{lemma}

\begin{proof}
This is the principal-three-dimensional / Coxeter-plane
decomposition of \cite{Kostant59}; the precise statement
needed here appears in \cite[Theorem~6.7 and the explicit
plane-wise rotation formula in \S 6.6--6.7,
pp.~983--988]{Kostant59}. The pairwise orthogonality of the
$P_k$ in the $W$-invariant inner product follows because the
real $c$-invariant decomposition is the realification of the
complex eigenspace decomposition of $c$ acting on $V \otimes
\C$, and the eigenvalue spectrum $\{\zeta^{m_k^\pm} : k = 1,
\ldots, 4\}$ (where $\zeta = e^{2\pi i / 30}$) consists of
distinct simple eigenvalues paired into complex conjugate
$\{\zeta^{m_k^-}, \zeta^{m_k^+}\} = \{\zeta^{m_k^-},
\overline{\zeta^{m_k^-}}\}$ pairs (each $m_k^+ = 30 - m_k^-$
gives the conjugate root of unity). For the eigenvalue
formula see also \cite[\S 3.20 Thm.~3.20 (the eigenvalues of
a Coxeter element)]{Humphreys}.
\end{proof}

\begin{lemma}[McKay correspondence: $2I \leftrightarrow E_8$]
\label{lem:McKay}
The binary icosahedral group $2I \subset SU(2)$ corresponds
to the affine Dynkin diagram $\widetilde{E}_8$ under the
McKay correspondence: the irreducible representations of
$2I$ are in canonical bijection with the nodes of
$\widetilde{E}_8$ (including the trivial representation on
the affine node), and the tensor-product graph
$V_{\rm std} \otimes (-)$ (where $V_{\rm std}$ is the
defining $2$-dimensional representation of $2I$)
reproduces the $\widetilde{E}_8$ adjacency.
\end{lemma}

\begin{proof}
\cite[\S 3, pp.~183--186]{McKay80}, where the ADE
correspondence between finite subgroups of $SU(2)$ and
extended Dynkin diagrams is established (the $2I
\leftrightarrow \widetilde E_8$ entry of the ADE table is
the load-bearing case for this paper).
\end{proof}

\begin{remark}[Role of Lemma~\ref{lem:McKay} in this paper]
\label{rmk:McKay-role}
Lemma~\ref{lem:McKay} is a contextual / language-bridge
statement only. We use it to justify the simultaneous use of
$E_8$ and $H_4$ / $600$-cell language for the same underlying
classical data, and to label the $\alpha$-1 phason complement
intrinsically by the $E_8$ Coxeter geometry. We do
\emph{not} extract the $600$-cell rational-irrep Laplacian
spectrum, the count $50$, or any $137$-arithmetic from
Lemma~\ref{lem:McKay}; those belong to $\alpha$-3 and to
separate downstream papers.
\end{remark}

%=====================================================================
\section{The Coxeter-Resolved Phason Module $\Mcox$}
\label{sec:Mcox}
%=====================================================================

We now introduce the Coxeter-resolved phason module $\Mcox$
and its reconciliation with $\alpha$-1's $\Mtw$.

\begin{definition}[The Coxeter-resolved phason module]
\label{def:Mcox}
Fix the four Coxeter planes $P_1, \ldots, P_4 \subset V =
\R^8$ of Lemma~\ref{lem:cox-plane}. The
\emph{Coxeter-resolved phason module} is the formal free
$\Ophi$-module
\[
  \Mcox \;:=\; \bigoplus_{k=1}^{4} \Ophi \cdot e_k,
\]
where $e_k$ is a labelling generator attached to plane $P_k$.
Thus $\rk_\Ophi \Mcox = 4$ and $\rk_\Z \Mcox = 8$.
\end{definition}

\begin{definition}[Reconciliation map; concrete construction]
\label{def:reconciliation}
Fix a Coxeter element $c \in W(E_8)$. For each $k \in \{1,
2, 3, 4\}$, let $u_k(c) \in P_k(c)$ be the
\emph{distinguished plane-wise generator} chosen by
Hypothesis $\Hor$ (Hypothesis~\ref{hyp:Horient}): the unit
real-vector along the upper-eigenline complex coordinate of
$c|_{P_k(c)}$ that $\Hor$ singles out (this fixes $u_k(c)$
up to sign, which is part of $\Hor$'s orientation choice).
The \emph{reconciliation map} is the $\sigma$-twisted
$\Ophi$-linear map
\[
  \Psi \colon \Mcox(c) \;\longrightarrow\; \HH_{\rm int}
\]
defined on the $\Ophi$-basis $\{e_1, e_2, e_3, e_4\}$ of
$\Mcox(c)$ by
\[
  \Psi(e_k) \;:=\; \pi_{\rm int}^{\rm amb}(u_k(c))
    \;\in\; \HH_{\rm int}, \qquad k = 1, 2, 3, 4,
\]
and extended $\sigma$-twisted-$\Ophi$-linearly:
$\Psi(\alpha \cdot e_k) = \sigma(\alpha) \cdot \Psi(e_k)$
for $\alpha \in \Ophi$. The codomain is the entire
$\HH_{\rm int}$ as a target $\R$-vector space; under
Hypothesis $\Hdist$ (Hypothesis~\ref{hyp:Hdist}), the image
$\Psi(\Mcox(c))$ is shown in
Proposition~\ref{prop:reconciliation} to land inside the
$\Ophi$-submodule $\pi_{\rm int}^{\rm amb}(\Mtw) \subseteq
\HH_{\rm int}$.
\end{definition}

\begin{hypothesis}[$\Hdist$: plane-wise distribution of the
  phason internal image]
\label{hyp:Hdist}
For the chosen Coxeter element $c$ and the orientation
convention $\Hor$, the four images $\{\Psi(e_1), \Psi(e_2),
\Psi(e_3), \Psi(e_4)\} \subset \pi_{\rm int}^{\rm
amb}(\Mtw)$ defined in Def.~\ref{def:reconciliation} are
non-zero and together generate the rank-$2$ $\Ophi$-submodule
$\pi_{\rm int}^{\rm amb}(\Mtw)$ of $\HH_{\rm int}$ as a
$\Z$-module. Geometrically: the $\sigma$-twisted internal
image of $\alpha$-1's $\Mtw$ distributes nontrivially across
all four Coxeter planes $P_k(c)$, each plane contributing a
non-zero $\Ophi$-line. \emph{We do not assert the four
$\Ophi$-lines are $\Ophi$-independent}: in a rank-$2$
$\Ophi$-module, four non-zero lines must satisfy at least two
$\Ophi$-linear dependencies. The dependency structure is left
unresolved in $\alpha$-2; the hypothesis controls only
generation, not linear independence of the four plane-image
lines.
\end{hypothesis}

\begin{remark}[Status of $\Hdist$]
\label{rmk:Hdist-status}
$\Hdist$ is a \emph{structural distribution hypothesis},
not a theorem of $\alpha$-2. Its content is that the
rank-$2$ $\Ophi$-image $\pi_{\rm int}^{\rm amb}(\Mtw)
\subset \HH_{\rm int}$ is generated, as a $\Z$-module, by four
nonzero plane-image lines (one per Coxeter plane), rather than
concentrating in only some subset of the planes. Since the
module has $\Ophi$-rank $2$, the four generating lines must
satisfy non-trivial $\Ophi$-dependencies; $\Hdist$ does not
fix the dependency structure. Geometrically: each $P_k(c)$ contributes
non-trivially to the internal phason image. A fully rigorous
derivation of $\Hdist$ from the bare $\alpha$-1 data would
require analysing the action of the centralizer of $c$ in
$W(E_8)$ on the phason internal image and proving
distribution-uniformity; we do not provide that derivation
here. The map $\Psi$ of Def.~\ref{def:reconciliation} is
already well-defined as a map $\Mcox(c) \to \HH_{\rm int}$
without $\Hdist$ (the formula $\Psi(e_k) = \pi_{\rm
int}^{\rm amb}(u_k(c))$ uses only $\Hor$). $\Hdist$ is
needed only to upgrade the codomain to $\pi_{\rm int}^{\rm
amb}(\Mtw)$ and to give the asserted image-rank, as in
Proposition~\ref{prop:reconciliation}.
\end{remark}

\begin{remark}[On the comparison level]
\label{rmk:reconciliation-level}
The reconciliation $\Psi$ is \emph{not} a module isomorphism
$\Mcox \not\cong \Mtw$ (they have different $\Z$-ranks: $8$
vs $4$). Rather, $\Psi$ resolves the internal-image data of
$\Mtw$ along the four orthogonal Coxeter planes: $\Mtw$
encodes the rank-$2$ $\Ophi$-module of internal phason data
attached to the cut-and-project, while $\Mcox$ records the
plane-wise contribution of each of the four Coxeter
directions to that data. The image of $\Psi$ has $\Ophi$-rank
$\leq 2$ inside $\HH_{\rm int}$ (consistent with $\alpha$-1's
admissibility axiom (A4)); the kernel of $\Psi$ records the
remaining $\Ophi$-rank $\geq 2$ of bookkeeping degeneracy.
The bookkeeping equivalence claim --- that $\Mcox$ and $\Mtw$
encode the same underlying internal geometry up to the
explicit map $\Psi$ --- is made precise in
Proposition~\ref{prop:reconciliation}.
\end{remark}

\begin{proposition}[Bookkeeping equivalence at the
  cut-and-project parameter level; conditional on $\Hmin
  + \Hor + \Hdist$]
\label{prop:reconciliation}
Fix a Coxeter element $c \in W(E_8)$ and assume all three:
\begin{itemize}
\item $\Hmin$, so that $\alpha$-1's Theorem $C_1$ applies
  and $\Mtw \cong \Ophi^{2}$ as an $\Ophi$-module;
\item Hypothesis $\Hor$ (Hypothesis~\ref{hyp:Horient}),
  fixing the plane-wise upper-eigenline coordinate
  $u_k(c) \in P_k(c)$;
\item Hypothesis $\Hdist$ (Hypothesis~\ref{hyp:Hdist}),
  ensuring the four images $\Psi(e_k) := \pi_{\rm int}^{\rm
  amb}(u_k(c))$ are non-zero and together generate
  $\pi_{\rm int}^{\rm amb}(\Mtw)$ as a $\Z$-module.
\end{itemize}
Then the reconciliation map $\Psi$ of
Definition~\ref{def:reconciliation} is well-defined,
$\sigma$-twisted-$\Ophi$-linear, and its image
$\Psi(\Mcox(c))$ has $\Ophi$-rank exactly equal to the
$\Ophi$-rank of $\pi_{\rm int}^{\rm amb}(\Mtw)$ (which is
$\leq 2$ by $\alpha$-1's (A4); by $\Hdist$, the four
$\Z$-generators span the full image). The kernel
$\ker(\Psi) \subseteq \Mcox(c)$ is the bookkeeping
degeneracy recording the redundancy between the four-plane
$\Mcox$ description and the rank-$2$ $\Mtw$ description.
\end{proposition}

\begin{proof}
\emph{Step 1: $\Psi$ is well-defined.} For each $k$, $u_k(c)
\in P_k(c) \subset V \otimes \R$ is a chosen unit vector
fixed by $\Hor$ (orientation choice). The image $\Psi(e_k) =
\pi_{\rm int}^{\rm amb}(u_k(c))$ is the $\R$-linear image
under the (extended) internal projection of $\alpha$-1's
(A3); since $\pi_{\rm int}^{\rm amb}$ is well-defined on
$V \otimes \R$, $\Psi(e_k)$ is a well-defined element of
$\HH_{\rm int}$. The extension to all of $\Mcox(c)$ via
$\sigma$-twisted-$\Ophi$-linearity is by the universal
property of free modules: any choice of basis-image
extends uniquely.

\emph{Step 2: $\Psi$ takes values in $\pi_{\rm int}^{\rm
amb}(\Mtw)$.} By $\Hdist$, $\Psi(e_k) \in \pi_{\rm
int}^{\rm amb}(\Mtw)$ for each $k$ (the hypothesis explicitly
locates the four images inside the phason image). Since
$\pi_{\rm int}^{\rm amb}(\Mtw)$ is an $\Ophi$-submodule of
$\HH_{\rm int}$ (under the $\sigma$-twisted action), the
$\sigma$-twisted-$\Ophi$-linear extension of $\Psi$ to all
of $\Mcox(c)$ takes values in this submodule.

\emph{Step 3: $\sigma$-twisted-$\Ophi$-linearity.}
$\Ophi$-linearity (with respect to $\sigma$-twisted
scalars) is by construction (the universal property of
$\Mcox(c)$ as a free $\Ophi$-module on $\{e_k\}$ combined
with the explicit twist convention). It is consistent with
$\alpha$-1 (A3.b): $\pi_{\rm int}^{\rm amb}$ is itself
$\sigma$-twisted $\Ophi$-linear, so the entire diagram
commutes.

\emph{Step 4: Image rank.}
$\Psi(\Mcox(c))$ is a finitely generated $\Ophi$-submodule
of $\pi_{\rm int}^{\rm amb}(\Mtw) \subseteq \HH_{\rm int}$,
hence (by $\alpha$-1's (A4)) has $\Ophi$-rank $\leq 2$. By
$\Hdist$, the four $\Z$-generators $\{\Psi(e_1), \ldots,
\Psi(e_4)\}$ generate $\pi_{\rm int}^{\rm amb}(\Mtw)$ as a
$\Z$-module; hence the $\Ophi$-image equals
$\pi_{\rm int}^{\rm amb}(\Mtw)$. The kernel is then
determined by the rank--nullity count: $\rk_\Ophi \Mcox(c)
- \rk_\Ophi \Psi(\Mcox(c)) = 4 - r$ where $r := \rk_\Ophi
\pi_{\rm int}^{\rm amb}(\Mtw) \in \{0, 1, 2\}$.
\end{proof}

\begin{remark}[$L_{16}$ is bookkeeping, not a closure theorem]
\label{rmk:L16-bookkeeping}
Define
\[
  L_{16} \;:=\; E_8 \oplus \Mcox
\]
formally. We do \emph{not} claim that $L_{16}$ satisfies an
analog of $\alpha$-1's Theorem $C_2$ (closure / termination at
$\Z$-rank 16); we make no statement about admissibility,
no uniqueness claim, and no termination claim about $L_{16}$.
The symbol $L_{16}$ is purely shorthand for the Coxeter-resolved
phason bookkeeping data $E_8 + \Mcox$ used in
$T_A$ and $T_B$. Any future paper attempting to elevate
$L_{16}$ to a closure status would need a distinct theorem
not contained here.
\end{remark}

%=====================================================================
\section{Theorem $T_A$: Phason--Coxeter Alignment}
\label{sec:TA}
%=====================================================================

\begin{theorem}[$\mathbf{T_A}$: phason--Coxeter alignment;
  conditional on $\Hmin + \Hor + \Hdist$ relative to a fixed
  Coxeter element]
\label{thm:TA}
Fix a Coxeter element $c \in W(E_8)$ as input data, and let
$\{P_1(c), P_2(c), P_3(c), P_4(c)\}$ be the four
pairwise-orthogonal $c$-invariant Coxeter planes of
Lemma~\ref{lem:cox-plane}. Let $\Mcox = \Mcox(c)$ be the
Coxeter-resolved phason module of Definition~\ref{def:Mcox}
attached to this choice (one $\Ophi$-line $\Ophi \cdot e_k$
labelled by $P_k(c)$). Assume all three:
\begin{itemize}
\item $\alpha$-1's Hypothesis $\Hmin$, so that Theorem $C_1$
  applies and $\Mtw \cong \Ophi^{2}$ as an $\Ophi$-module;
\item Hypothesis $\Hor$ (Hypothesis~\ref{hyp:Horient}), the
  $\sigma$-twist orientation convention on Coxeter planes
  fixing the plane-by-plane upper-eigenline coordinate
  $u_k(c) \in P_k(c)$;
\item Hypothesis $\Hdist$ (Hypothesis~\ref{hyp:Hdist}), the
  distribution hypothesis ensuring the four images
  $\Psi(e_k) = \pi_{\rm int}^{\rm amb}(u_k(c))$ are non-zero
  and together generate $\pi_{\rm int}^{\rm amb}(\Mtw)$ as
  $\Z$-module.
\end{itemize}
Then:
\begin{enumerate}
\item[\textup{(a)}] The reconciliation map $\Psi$ of
  Definition~\ref{def:reconciliation} is well-defined and
  $\sigma$-twisted-$\Ophi$-linear, $\Psi \colon \Mcox(c) \to
  \pi_{\rm int}^{\rm amb}(\Mtw)$, with image of $\Ophi$-rank
  $\leq 2$ inside $\HH_{\rm int}$ (matching $\alpha$-1's
  bound (A4)).
\item[\textup{(b)}] $\Mcox(c) = \bigoplus_{k=1}^{4} \Ophi
  \cdot e_k$ as $\Ophi$-modules, with the four $\Ophi$-lines
  in bijection with $\{P_1(c), \ldots, P_4(c)\}$.
\item[\textup{(c)}] The bijection in (b) is canonical
  \emph{relative to the chosen Coxeter element $c$}, up to
  signed permutation $S_4 \ltimes (\Z/2)^4$ of the
  four-plane labels (ordering + orientation).
\end{enumerate}
\end{theorem}

\begin{proof}
\emph{Step 1: Reconciliation map (Part (a)).}
This is the content of Proposition~\ref{prop:reconciliation}
under the joint hypothesis $\Hmin + \Hor + \Hdist$. The
construction of $\Psi$ is plane-by-plane via the explicit
formula $\Psi(e_k) = \pi_{\rm int}^{\rm amb}(u_k(c))$ of
Def.~\ref{def:reconciliation}: $u_k(c)$ is fixed by $\Hor$,
the image lies in $\pi_{\rm int}^{\rm amb}(\Mtw)$ by
$\Hdist$, and the $\sigma$-twisted-$\Ophi$-linearity follows
from $\alpha$-1's (A3.b) on $\pi_{\rm int}^{\rm amb}$. The
image-rank bound is from $\alpha$-1's (A4); equality with
$\pi_{\rm int}^{\rm amb}(\Mtw)$'s $\Ophi$-rank uses
$\Hdist$ (the four $\Z$-generators span the full image).

\emph{Step 2: Direct-sum structure (Part (b)).}
By Definition~\ref{def:Mcox}, $\Mcox(c)$ is the formal free
$\Ophi$-module on the four generators $\{e_1, e_2, e_3,
e_4\}$, each labelled by one of the four $c$-invariant
Coxeter planes $\{P_1(c), \ldots, P_4(c)\}$. The direct-sum
decomposition $\Mcox(c) = \bigoplus_{k=1}^{4} \Ophi \cdot
e_k$ is therefore tautological from the construction. The
bijection between the $\Ophi$-lines $\{\Ophi \cdot e_k\}$
and the planes $\{P_k(c)\}$ is the labelling fixed in
Def.~\ref{def:Mcox}.

\emph{Step 3: Canonicity (Part (c)).}
By Lemma~\ref{lem:cox-plane}, the four $c$-invariant
pairwise-orthogonal planes $\{P_1(c), \ldots, P_4(c)\}$ are
determined by the chosen Coxeter element $c$ as an unordered
set. Once $c$ is fixed, the only further choices are:
\begin{enumerate}
\item[(i)] An ordering of the four planes: an $S_4$
  relabelling of $\{1, 2, 3, 4\} \to \{P_1(c), \ldots,
  P_4(c)\}$.
\item[(ii)] An orientation on each plane: an independent
  $\Z/2$ choice per plane, corresponding to the choice of
  generator $e_k$ vs $-e_k$ inside the $\Ophi$-line $\Ophi
  \cdot e_k$.
\end{enumerate}
Together these give an $S_4 \ltimes (\Z/2)^4$ action on
the labelling. The decomposition is rigid beyond this
signed-permutation quotient, conditional on $c$.

\emph{Note on different choices of $c$:} For two distinct
Coxeter elements $c, c' \in W(E_8)$, the planes
$\{P_k(c)\}$ and $\{P_k(c')\}$ are \emph{conjugate} in
$W(E_8)$ (since all Coxeter elements form a single
conjugacy class in $W(E_8)$, see
\cite[\S 3.18]{Humphreys}) but are generally not literally
equal as subspaces of $V$; we make no canonicity claim
across distinct choices of $c$.
\end{proof}

\begin{remark}[On the meaning of canonicity here]
\label{rmk:TA-canonicity}
Theorem~\ref{thm:TA} does \emph{not} claim that the four
specific $\Ophi$-lines of $\Mcox$ are absolutely intrinsic to
$E_8$. The chosen Coxeter element $c \in W(E_8)$ is part of
the input data of the theorem; different choices of $c$
give conjugate (but generally not identical) plane
decompositions of $V$. Relative to a fixed $c$, the
four-plane decomposition of $\Mcox(c)$ is canonical up to
the $S_4 \ltimes (\Z/2)^4$ signed-permutation action on
labels; beyond that, it is rigid. Across different choices
of $c$, the canonical content is only that the four-plane
decomposition exists and is unique up to conjugacy.
\end{remark}

%=====================================================================
\section{$\sigma$-Twist Orientation and Plane-Wise Winding}
\label{sec:sigma-twist}
%=====================================================================

We now make precise the $\sigma$-twist orientation reversal
and define the plane-wise winding integer.

\begin{lemma}[$\sigma$ commutes with $c$]
\label{lem:sigma-c-commute}
Let $c \in W(E_8)$ act on $V \otimes \Q$ via the rational
reflection representation. The Galois involution
$\sigma \in \Gal(K/\Q)$ extends entrywise to $\Ophi$-coefficient
vectors and acts on $V \otimes K$ by acting on coefficients.
Since $c$ has matrix entries in $\Z \subset \Ophi^\sigma$ (the
fixed subring of $\sigma$), $c$ commutes with $\sigma$ on
$V \otimes K$:
\[
  \sigma(c \cdot v) \;=\; c \cdot \sigma(v)
  \qquad \text{for all } v \in V \otimes K.
\]
\end{lemma}

\begin{proof}
$c$ is an element of $W(E_8)$, which acts on the integer root
lattice $E_8 \subset V$ with integer matrix entries
(\cite[Ch.~6, \S1]{Bourbaki46} or \cite[\S 1]{Humphreys}). On
$V \otimes K$, $c$'s matrix lies in $\mathrm{GL}_8(\Z)
\subset \mathrm{GL}_8(\Ophi)$. The Galois action $\sigma$
fixes elements of $\Z$ pointwise. Hence for any matrix
entry $c_{ij} \in \Z$, $\sigma(c_{ij}) = c_{ij}$, and the
matrix $c$ commutes with $\sigma$ on coefficients:
$\sigma(c \cdot v) = \sigma(c) \cdot \sigma(v) = c \cdot
\sigma(v)$.
\end{proof}

\begin{remark}[Why coefficientwise $\sigma$ does not act on
  the $c$-eigenspaces]
\label{rmk:sigma-no-action}
Although Lemma~\ref{lem:sigma-c-commute} gives $\sigma c =
c \sigma$ on $V \otimes K$, this does \emph{not} imply that
$\sigma$ acts non-trivially on the complex eigenspaces of
$c$ on $V \otimes \C$. The complex eigenvalues of $c$ are
$\zeta^{m_k^\pm}$ with $\zeta = e^{2\pi i / 30}$ and live
in the cyclotomic field $\Q(\zeta) \supset K$; complex
conjugation $\zeta \mapsto \zeta^{-1}$ on $\Q(\zeta)$ fixes
the totally real subfield $K = \Q(\sqrt 5)$ pointwise, so
it does \emph{not} restrict to the nontrivial $\sigma \in
\Gal(K/\Q)$. Lifts of $\sigma$ to $\Gal(\Q(\zeta)/\Q)$
exist (since $[\Q(\zeta) : K] = 4$), but each requires an
additional choice. In particular, the heuristic
``$\sigma$ on $V$ swaps the complex eigenspaces of $c$
within each $P_k$'' that appears in informal working notes
\emph{cannot} be derived from coefficientwise Galois
action on $V \otimes K$ alone, and as a global identity on
$V \otimes \C$ would contradict the commutation $\sigma c
= c \sigma$ (since $\sigma c \sigma = c$, not $c^{-1}$).

The orientation reversal needed for the upper-exponent
winding count below is therefore not an extension of
$\sigma$ to $V \otimes \C$. It is a separate structural
hypothesis $\Hor$ on the choice of plane-wise complex
identification adapted to $\alpha$-1's $\sigma$-twisted
internal projection (A3.b), formalised in
Hypothesis~\ref{hyp:Horient} below.
\end{remark}

\begin{hypothesis}[$\Hor$: $\sigma$-twist orientation
  convention on Coxeter planes]
\label{hyp:Horient}
Let $\Pi^{\rm amb} = (\pi_{\rm phys}^{\rm amb}, \pi_{\rm
int}^{\rm amb})$ be $\alpha$-1's admissible projection pair,
with $\pi_{\rm int}^{\rm amb}$ the $\sigma$-twisted
$\Ophi$-linear internal projection of $\alpha$-1's axiom
(A3.b) \cite[(A3.b)]{Cascade12D}. The hypothesis $\Hor$
asserts: for each Coxeter plane $P_k$
(Lemma~\ref{lem:cox-plane}), the plane-wise complex
identification $P_k \cong \C$ adapted to the
$\sigma$-twisted internal projection $\pi_{\rm int}^{\rm
amb}$ uses the \emph{upper} eigenline of $c$ (eigenvalue
$\zeta^{m_k^+}$) for the canonical complex coordinate,
rather than the lower eigenline (eigenvalue $\zeta^{m_k^-}$)
which is the natural $V$-side choice. Equivalently: under
$\Hor$, the plane-wise rotation generator $c|_{P_k}$
realises in the internal-factor identification as
$\R$-linear rotation by angle $2\pi m_k^+ / 30$ (rather than
by $2\pi m_k^- / 30$ as on the $V$-side).
\end{hypothesis}

\begin{remark}[Status of $\Hor$ and motivation]
\label{rmk:Horient-status}
$\Hor$ is a \emph{structural orientation hypothesis}, not a
theorem of this paper. Its motivation is that the
$\sigma$-twisted internal projection of $\alpha$-1's (A3.b)
treats the internal factor as the algebraic
``Galois-conjugate'' picture relative to the physical
factor: $\pi_{\rm phys}^{\rm amb}$ uses the standard
$\Ophi$-action, while $\pi_{\rm int}^{\rm amb}$ uses the
$\sigma$-twisted action $\alpha \cdot w := \sigma(\alpha)
\cdot w$. Inside each Coxeter plane $P_k$, the two complex
eigenlines of $c$ are interchanged by complex conjugation
of the eigenvalue (an $\R$-orientation reversal of the
plane). $\Hor$ is the assertion that the $\sigma$-twist of
(A3.b), \emph{interpreted plane-by-plane}, selects the
upper-eigenline complex coordinate, hence the upper-exponent
rotation labelling, on each $P_k$.

A fully rigorous derivation of $\Hor$ from the bare
algebraic data of (A3.b) would require formalising a
$\C$-structure on $\HH_{\rm int}$ compatible with the
$\sigma$-twisted $\Ophi$-action and proving canonicity of
the upper-eigenline choice plane-by-plane. We do not provide
that derivation here; introducing $\Hor$ as a named
hypothesis makes this orientation convention an explicit
input to $T_B$ rather than a hidden assumption.
\end{remark}

\begin{definition}[Plane-wise phason winding integer]
\label{def:winding}
For each Coxeter plane $P_k$ (Lemma~\ref{lem:cox-plane}), the
\emph{phason winding integer} $w_k$ is the unique integer
in $\{1, 2, \ldots, 29\}$ such that the plane-wise rotation
generator $c|_{P_k}$ acts in the $\sigma$-twist orientation
assumed by Hypothesis $\Hor$ as $\R$-linear rotation of the
internal factor by angle $2\pi w_k / 30$. Definition of $w_k$
is independent of the assertion that $w_k = m_k^+$; identifying
$w_k$ with a specific Coxeter exponent is the content of
Hypothesis $\Hor$, not of this definition.
\end{definition}

%=====================================================================
\section{Theorem $T_B$ (Total Winding $= 88$)}
\label{sec:TB-88}
%=====================================================================

\begin{proposition}[$\mathbf{T_B}$, total winding under the hypothesis stack]
\label{thm:TB-88}
\emph{The content of this statement is almost entirely packed
into the hypotheses $\Hor$ and $\Hdist$; we present it as a
named proposition to organise the eventual bookkeeping, not as
a substantive independent derivation.} Assume all three:
\begin{itemize}
\item $\alpha$-1's Hypothesis $\Hmin$
  (\cite{Cascade12D}), so that Theorem~$C_1$ applies;
\item Hypothesis~\ref{hyp:Horient} ($\Hor$, the
  $\sigma$-twist orientation convention assigning each plane-wise
  winding integer $w_k$ to the upper exponent $m_k^+$);
\item Hypothesis~\ref{hyp:Hdist} ($\Hdist$, the plane-wise
  distribution hypothesis), so that Theorem $T_A$
  (Theorem~\ref{thm:TA}) holds for the chosen Coxeter element
  $c$.
\end{itemize}
Then each Coxeter plane $P_k(c)$ has phason winding integer
$w_k = m_k^+$ (\emph{by the content of} $\Hor$), and the total
phason winding is the Coxeter-exponent sum
\[
  \Wph \;:=\; \sum_{k=1}^{4} w_k
       \;=\; \sum_{k=1}^{4} m_k^+
       \;=\; 17 + 19 + 23 + 29
       \;=\; 88,
\]
the final equality being a classical $E_8$ exponent-arithmetic
identity (Lemma~\ref{lem:E8-data}).
\end{proposition}

\begin{proof}
\emph{Step 1: $T_A$ provides the four-plane phason
decomposition under the hypothesis stack.}
Theorem~\ref{thm:TA} (under $\Hmin + \Hor + \Hdist$) provides
the canonical four-plane decomposition of the Coxeter-resolved
phason bookkeeping module: one $\Ophi$-line per Coxeter plane
$P_k$, with each $\Ophi$-line carrying the $\sigma$-twisted
internal-factor data of plane $P_k$.

\emph{Step 2: Under $\Hor$, each $w_k$ is identified with $m_k^+$.}
Hypothesis $\Hor$ asserts exactly this identification; we do not
derive it independently.

\emph{Step 3: Sum.} By Lemma~\ref{lem:E8-data}, the upper
exponent set of $E_8$ is $\{17, 19, 23, 29\}$, and
\[
  \sum_{k=1}^{4} m_k^+ \;=\; 17 + 19 + 23 + 29 \;=\; 88.
\]
This last step is the only independent mathematical content
beyond the hypothesis stack.
\end{proof}

\begin{remark}[Conditionality fence for Theorem~\ref{thm:TB-88}]
\label{rmk:TB-88-fence}
Theorem~\ref{thm:TB-88} is conditional on the conjunction
$\Hmin + \Hor + \Hdist$ (the $\alpha$-1 minimality
hypothesis via $T_A$, the $\sigma$-twist orientation
convention, and the plane-wise distribution hypothesis). It
is \emph{unconditional} in $\Hquot$: no quotient hypothesis
is used here. The total $88$ is an honest theorem at the
$\Hmin + \Hor + \Hdist$ conditional level.
\end{remark}

%=====================================================================
\section{Theorem $T_B$ (Slot Count $= 87$ under $\Hquot$)}
\label{sec:TB-87}
%=====================================================================

The slot-count clause $\Nslot = 87$ requires identifying and
quotienting out a one-dimensional global-phase artifact. We
introduce the quotient hypothesis explicitly.

\begin{definition}[Candidate global-phase quotient direction]
\label{def:global-phase}
Among the $88$ phason winding units of $\Wph$
(Proposition~\ref{thm:TB-88}), there is a natural candidate
\emph{common-winding-mode direction}: the combinatorial diagonal
$\{(w_1, w_2, w_3, w_4) : w_1 = w_2 = w_3 = w_4 \bmod 30\}$ of
the four-plane phason grading under the $\Z/30\Z$-order of each
Coxeter-plane rotation $c|_{P_k}$. This is a candidate diagonal
mode in the bookkeeping grading, \emph{not} an independently
constructed $U(1)$-action. A genuine $U(1)$-action realising
this direction --- and in particular fixing its contribution to
the winding count --- requires an additional construction not
carried out here. We do not assert that this candidate direction
has dimension exactly $1$ inside $\Wph$, nor that quotienting by
it removes exactly one winding unit; both claims are content of
Hypothesis~\ref{hyp:Hquot} below.
\end{definition}

\begin{remark}[The need for a separate quotient hypothesis]
\label{rmk:Hquot-justif}
Definition~\ref{def:global-phase} identifies only a
\emph{candidate} quotient direction, not its dimension or
its winding contribution. A fully rigorous derivation that
quotienting by this direction removes exactly $1$ unit of
winding (and not, say, $0$ or $4$) requires a precise
quotient construction: explicitly identifying a global
$U(1)$-symmetry action on the Coxeter-resolved phason
bookkeeping and proving that its quotient on $\Wph$ has
the asserted dimension.

Such a construction is well beyond the scope of $\alpha$-2;
it would require introducing the structure of the
cut-and-project's acceptance window as a $U(1)$-bundle over
the Coxeter base. We therefore proceed by introducing the
quotient hypothesis $\Hquot$ explicitly, as the assumption
that the candidate direction of
Def.~\ref{def:global-phase} contributes exactly one
winding unit to $\Wph$:
\end{remark}

\begin{hypothesis}[$\Hquot$: one-dimensional global-phase
  quotient]
\label{hyp:Hquot}
The $88$ units of phason winding accumulated in
$\Wph$ (Theorem~\ref{thm:TB-88}) admit a quotient by exactly
one dimension of global rotational phase
(Def.~\ref{def:global-phase}), giving a slot count
$\Nslot = \Wph - 1 = 87$.
\end{hypothesis}

\begin{proposition}[$\mathbf{T_B}$, slot count, restatement under the hypothesis stack]
\label{thm:TB-87}
\emph{The arithmetic content is $88 - 1 = 87$; the mathematical
content is packed into Hypothesis $\Hquot$. Recorded as a named
proposition for bookkeeping, not as a substantive independent
derivation.} Assume all four:
\begin{itemize}
\item $\alpha$-1's Hypothesis $\Hmin$;
\item Hypothesis $\Hor$ (Hypothesis~\ref{hyp:Horient});
\item Hypothesis $\Hdist$ (Hypothesis~\ref{hyp:Hdist}); so
  that $T_A$ Theorem~\ref{thm:TA} and Theorem~\ref{thm:TB-88}
  give $\Wph = 88$;
\item Hypothesis $\Hquot$ (Hypothesis~\ref{hyp:Hquot}, the
  one-dimensional global-phase quotient).
\end{itemize}
Then the phason slot count is
\[
  \Nslot \;=\; \Wph - 1 \;=\; 88 - 1 \;=\; 87.
\]
\end{proposition}

\begin{proof}
By Proposition~\ref{thm:TB-88}, $\Wph = 88$. By $\Hquot$ (which
directly asserts the quotient removes exactly one winding unit),
$\Nslot = \Wph - 1 = 87$. The arithmetic is direct; no additional
mathematical content beyond the hypothesis stack.
\end{proof}

\begin{remark}[Status of Theorem~\ref{thm:TB-87}]
\label{rmk:TB-87-status}
Theorem~\ref{thm:TB-87} is conditional on the conjunction
$\Hmin + \Hor + \Hdist + \Hquot$: the $\alpha$-1 minimality
hypothesis, the $\sigma$-twist orientation convention, the
plane-wise distribution hypothesis, and the one-dimensional
global-phase quotient. We do
\emph{not} claim any of these four hypotheses is itself a
theorem of this paper. They are named hypotheses with the
structural justifications of Remarks~\ref{rmk:Horient-status},
\ref{rmk:Hdist-status}, \ref{rmk:Hquot-justif}.
A future paper that constructs the explicit
$U(1)$-bundle quotient on the cut-and-project acceptance
window would close $\Hquot$ to a theorem; in the meantime,
the slot count $87$ is conditionally derived. The total
$\Wph = 88$ stands at the $T_A$-conditional level without
$\Hquot$.
\end{remark}

%=====================================================================
\section{Self-Audit}
\label{sec:audit}
%=====================================================================

\textbf{Unconditional layer.}
\begin{itemize}
\item[[1]] Lemma~\ref{lem:E8-data} ($E_8$ Coxeter data) cites
  classical sources at theorem/page level
  (Bourbaki Ch.~6 \S4 No.~12, Planche~VII; Humphreys
  \S 3.18--3.20).
\item[[2]] Lemma~\ref{lem:cox-plane} (Coxeter-plane
  decomposition) cites Kostant 1959 Theorem 6.7 with
  page-level pin.
\item[[3]] Lemma~\ref{lem:McKay} (McKay correspondence) cites
  McKay 1980 \S3, pp.~183--186.
\item[[4]] Lemma~\ref{lem:sigma-c-commute} ($\sigma$-$c$
  commutation) is a direct calculation from the
  $\Z$-integrality of $c$'s matrix.
\end{itemize}

\textbf{Structural hypotheses (named, not claimed as
theorems of $\alpha$-2).}
\begin{itemize}
\item[[5]] Hypothesis~\ref{hyp:Horient} ($\Hor$, the
  $\sigma$-twist orientation convention on Coxeter planes)
  is a \emph{named structural hypothesis}, not a proved
  matrix lemma; coefficient-Galois alone cannot derive
  plane-wise orientation reversal
  (Remark~\ref{rmk:sigma-no-action}). Structural motivation
  in Remark~\ref{rmk:Horient-status}.
\item[[5b]] Hypothesis~\ref{hyp:Hdist} ($\Hdist$, the
  plane-wise distribution hypothesis on the phason internal
  image) is a \emph{named structural hypothesis}, not a
  derived theorem. Structural motivation in
  Remark~\ref{rmk:Hdist-status}: a full derivation would
  require analysing the centralizer-of-$c$ action on the
  phason internal image.
\end{itemize}

\textbf{Conditional on $\Hmin + \Hor + \Hdist$.}
\begin{itemize}
\item[[6]] $\Hmin$ is imported from $\alpha$-1
  (Fact~\ref{fact:alpha1-imports}); not claimed as a
  theorem of $\alpha$-2.
\item[[7]] Theorem~\ref{thm:TA} ($T_A$: phason--Coxeter
  alignment) conditional on $\Hmin + \Hor + \Hdist$,
  relative to a fixed Coxeter element $c \in W(E_8)$. The
  reconciliation map $\Psi$ is constructed explicitly
  plane-by-plane in Def.~\ref{def:reconciliation} (using
  $u_k(c)$ from $\Hor$ and image-distribution from
  $\Hdist$). Canonicity class explicitly stated as ``up to
  signed permutation of the four plane labels relative to
  $c$'' (Remark~\ref{rmk:TA-canonicity}); not overclaimed
  as rigid or canonical across distinct $c$.
\item[[8]] Theorem~\ref{thm:TB-88} (total winding $\Wph =
  88$) conditional on $\Hmin + \Hor + \Hdist$ via $T_A$
  together with the plane-wise upper-exponent labelling of
  $\Hor$. Independent of $\Hquot$.
\end{itemize}

\textbf{Conditional on $\Hmin + \Hor + \Hdist + \Hquot$.}
\begin{itemize}
\item[[9]] Hypothesis $\Hquot$ stated explicitly as
  Hypothesis~\ref{hyp:Hquot}; structurally justified in
  Remark~\ref{rmk:Hquot-justif}; not claimed as a theorem.
\item[[10]] Theorem~\ref{thm:TB-87} (slot count $\Nslot =
  87$) conditional on $\Hmin + \Hor + \Hdist + \Hquot$. The quotient
  step is identified as the residual gap.
\end{itemize}

\textbf{Convention discipline.}
\begin{itemize}
\item[[11]] The rank convention (this paper uses
  $\rk_\Ophi$ as default; $\alpha$-1 used $\rk_\Z$ as
  default) is fixed once in
  Notation~\ref{not:rank-conv} and the bridge to $\alpha$-1
  is made explicit in \S\ref{sec:conv-recon} and
  Proposition~\ref{prop:reconciliation}.
\item[[12]] $L_{16} = E_8 \oplus \Mcox$ is bookkeeping
  (Remark~\ref{rmk:L16-bookkeeping}); no closure-theorem
  claim attached.
\item[[13]] Reconciliation map $\Psi$ is defined explicitly
  (Def.~\ref{def:reconciliation}); image $\Ophi$-rank
  $\leq 2$ matches $\alpha$-1's (A4) bound
  (Prop.~\ref{prop:reconciliation}).
\end{itemize}

\textbf{Not claimed.}
\begin{itemize}
\item No unconditional $T_A$ or $T_B$ (slot-count clause).
\item No claim that $L_{16}$ replaces $L_{12}$ as the
  canonical upward closure.
\item No $\alpha^{-1} = 137 + \pi/87$ derivation.
\item No $137 = 87 + 50$ decomposition.
\item No $600$-cell Laplacian spectrum.
\item No photon, microtubule, biology, Standard-Model claim.
\item No dependence on P1--P7.
\item No use of \texttt{cascade-derivation/} working notes
  as proof source.
\end{itemize}

%=====================================================================
\section{Citation Contract and Downstream Handoff}
\label{sec:handoff}
%=====================================================================

\textbf{Unconditionally citable.}
\begin{itemize}
\item Lemmas~\ref{lem:E8-data}, \ref{lem:cox-plane},
  \ref{lem:McKay} (classical $E_8$ / Coxeter / McKay data).
\item Lemma~\ref{lem:sigma-c-commute} ($\sigma$ commutes with
  $c$ via $\Z$-integrality of $c$).
\end{itemize}

\textbf{Conditional on $\Hmin + \Hor + \Hdist$.}
\begin{itemize}
\item Theorem~\ref{thm:TA} ($T_A$: four-plane alignment of
  $\Mcox$, relative to a fixed Coxeter element $c$).
\item Theorem~\ref{thm:TB-88} (total winding $\Wph = 88$),
  via $T_A$ together with the plane-wise upper-exponent
  labelling of $\Hor$.
\end{itemize}

\textbf{Conditional on $\Hmin + \Hor + \Hdist + \Hquot$.}
\begin{itemize}
\item Theorem~\ref{thm:TB-87} (slot count $\Nslot = 87$).
\end{itemize}

\textbf{Downstream usage ($\alpha$-3).} The $\alpha$-3 paper
(fine-structure synthesis $\alpha^{-1} = 137 + \pi/87$) may
cite Theorem~\ref{thm:TA} and Theorem~\ref{thm:TB-88} with
their $\Hmin + \Hor + \Hdist$ conditional fence, and
Theorem~\ref{thm:TB-87} with the additional $\Hquot$
conditional fence. The $50$ count (rational-irrep Laplacian
on the $600$-cell) and the $137 = 87 + 50$ arithmetic are
\emph{not} provided here and must be supplied separately
by $\alpha$-3.

\textbf{Does NOT close.}
\begin{itemize}
\item $\Hmin$ as a theorem (deferred to a follow-up of
  $\alpha$-1).
\item $\Hor$ as a theorem: the $\sigma$-twist plane-wise
  orientation convention is left as a named structural
  hypothesis (Hypothesis~\ref{hyp:Horient}); a full derivation
  would require formalising a canonical $\C$-structure on
  $\HH_{\rm int}$ compatible with $\alpha$-1's (A3.b)
  $\sigma$-twisted action, which is beyond the scope of
  $\alpha$-2 (see Remark~\ref{rmk:Horient-status}).
\item $\Hdist$ as a theorem: the distribution hypothesis on
  the phason internal image (Hypothesis~\ref{hyp:Hdist})
  is a named structural input; a full derivation would
  require analysing the centralizer-of-$c$ action on
  $\pi_{\rm int}^{\rm amb}(\Mtw)$ and proving uniform
  plane-wise distribution, which is beyond the scope of
  $\alpha$-2 (see Remark~\ref{rmk:Hdist-status}).
\item $\Hquot$ as a theorem: the explicit one-dimensional
  global-phase quotient construction (a $U(1)$-bundle
  structure on the cut-and-project acceptance window) is
  beyond the scope of $\alpha$-2.
\item Any $\alpha^{-1}$-level fine-structure claim.
\end{itemize}

%=====================================================================
\section{Notation Audit}
\label{sec:notation}
%=====================================================================

\begin{center}
\begin{tabular}{|l|l|l|}
\hline
Symbol & Meaning & First use \\
\hline
$\varphi$ & $(1+\sqrt 5)/2$, positive fixed point of
  $\mathbf{F1}$ & Axiom~\ref{ax:F1} \\
$K, \Ophi, \sigma$ & $\Q(\sqrt 5), \Z[\varphi], $ Galois
  involution & Notation~\ref{not:K} \\
$h$ & Coxeter number of $E_8$, $h = 30$ &
  Lemma~\ref{lem:E8-data} \\
$c$ & Coxeter element, $c \in W(E_8)$ &
  Lemma~\ref{lem:E8-data} \\
$m_k^-$, $m_k^+$ & lower / upper exponent of conjugate pair
  attached to $P_k$ & Lemma~\ref{lem:E8-data} \\
$P_k$ ($k=1,\ldots,4$) & Coxeter plane in $V = \R^8$ &
  Lemma~\ref{lem:cox-plane} \\
$\Mtw$ & $\alpha$-1 phason complement, $\rk_\Z = 4$,
  $\rk_\Ophi = 2$ & Fact~\ref{fact:alpha1-imports} \\
$\Mcox$ & Coxeter-resolved phason module, $\rk_\Ophi = 4$,
  $\rk_\Z = 8$ & Def.~\ref{def:Mcox} \\
$\Psi$ & Reconciliation map $\Mcox(c) \to \HH_{\rm int}$
  (lands in $\pi_{\rm int}^{\rm amb}(\Mtw)$ under $\Hdist$)
  & Def.~\ref{def:reconciliation}, Prop.~\ref{prop:reconciliation} \\
$L_{16}$ & $E_8 \oplus \Mcox$, bookkeeping only &
  Rem.~\ref{rmk:L16-bookkeeping} \\
$w_k$ & plane-wise winding integer on $P_k$ &
  Def.~\ref{def:winding} \\
$\Wph$ & total phason winding, $\sum_k w_k = 88$ &
  Thm.~\ref{thm:TB-88} \\
$\Nslot$ & phason slot count after $\Hquot$ quotient, $= 87$ &
  Thm.~\ref{thm:TB-87} \\
$\Hmin$ & $\alpha$-1 minimality hypothesis on phason ring &
  Fact~\ref{fact:alpha1-imports} \\
$\Hor$ & $\sigma$-twist orientation convention on Coxeter
  planes & Hyp.~\ref{hyp:Horient} \\
$\Hdist$ & plane-wise distribution hypothesis on phason
  internal image & Hyp.~\ref{hyp:Hdist} \\
$\Hquot$ & one-dimensional global-phase quotient hypothesis &
  Hyp.~\ref{hyp:Hquot} \\
$T_A, T_B$ & Theorem~\ref{thm:TA}; Theorems~\ref{thm:TB-88},
  \ref{thm:TB-87} & \\
\hline
\end{tabular}
\end{center}

\begin{thebibliography}{99}

\bibitem{CascadeFoundations}
\emph{Cascade Correspondence Foundations},
\texttt{papers/cascade-correspondence-foundations/foundations.tex}.
Cited solely for the statement of Axiom $\mathbf{F1}$ (\S F1).
Not cited as a theorem-grade proof source.

\bibitem{Cascade12D}
\emph{The 12D Closure Theorem: Phason Complement Uniqueness
and Self-Reference on the Upward Cascade} ($\alpha$-1),
\texttt{papers/cascade-12d-closure/cascade-12d-closure.tex}.
Theorems $C_1$, $C_2$, $C_3$, Hypothesis $\Hmin$,
admissibility axioms (A1)--(A4), and the
$\sigma$-twist convention (A3.b).

\bibitem{Bourbaki46}
N.~Bourbaki, \emph{Groupes et alg\`ebres de Lie}, Chapitres
4--6. Hermann, Paris, 1968. (Reprinted Springer, Berlin,
2002.) For $E_8$: Ch.~6, \S4 No.~12 and Planche~VII (the
$E_8$ exponent table). Coxeter elements and exponent
symmetry: Ch.~5, \S6.

\bibitem{Humphreys}
J.~E.~Humphreys, \emph{Reflection Groups and Coxeter
Groups}, Cambridge Studies in Advanced Mathematics 29,
Cambridge University Press, 1990. \S 3.7 (table of
exponents for irreducible Coxeter groups including $E_8$);
\S 3.16 (Coxeter element); \S 3.18 (Coxeter number);
\S 3.19--3.20 (eigenvalues / exponents of a Coxeter
element).

\bibitem{Kostant59}
B.~Kostant, \emph{The principal three-dimensional subgroup
and the Betti numbers of a complex simple Lie group}, Amer.
J. Math. \textbf{81} (1959), 973--1032. Theorem~6.7
and \S 6.6--6.7 (pp.~983--988): the
Coxeter-plane decomposition of the reflection
representation into pairwise-orthogonal $c$-invariant
$2$-planes with explicit rotation angles.

\bibitem{ConwaySloane}
J.~H.~Conway, N.~J.~A.~Sloane, \emph{Sphere Packings,
Lattices and Groups}, 3rd ed., Springer, 1999. The icosian
ring $I$ and the $E_8$ embedding via icosians: Ch.~8,
\S2 (pp.~206--209), in particular \S 2.1 (icosian ring) and
\S 2.2 (the $E_8$ embedding via $(x, \sigma(x))$); table 8.1
on p.~209 lists the $240$ minimal vectors.

\bibitem{MoodyPatera}
R.~V.~Moody, J.~Patera, \emph{Quasicrystals and icosians},
J.~Phys.~A \textbf{26} (1993), 2829--2853. Theorem~4.1 and
Proposition~4.2: the explicit $E_8 \to H_4$ cut-and-project
interface used in $\alpha$-1's classical package.

\bibitem{ElserSloane}
V.~Elser, N.~J.~A.~Sloane, \emph{A highly symmetric
four-dimensional quasicrystal}, J.~Phys.~A \textbf{20}
(1987), 6161--6167. Original cut-and-project quasicrystal
construction from $E_8$: \S 2 (pp.~6162--6163) sets up the
$I \oplus \sigma(I)$ decomposition; \S 3 (pp.~6163--6165)
defines the acceptance window and projects onto
$\HH_{\rm phys}$.

\bibitem{McKay80}
J.~McKay, \emph{Graphs, singularities, and finite groups},
in: \emph{The Santa Cruz Conference on Finite Groups},
Proc.~Symp.~Pure Math.~37, Amer.~Math.~Soc., 1980,
pp.~183--186. \S 3 establishes the ADE correspondence
between finite subgroups of $SU(2)$ and extended Dynkin
diagrams; the $2I \leftrightarrow \widetilde E_8$ entry is
in the body of this section.

\bibitem{Marcus}
D.~A.~Marcus, \emph{Number Fields}, 2nd ed., Universitext,
Springer, 2018. Ring of integers of $\Q(\sqrt d)$ for
$d \equiv 1 \pmod 4$: Ch.~2, Thm.~1.

\end{thebibliography}

\end{document}
